\section{Cumulative distribution functions}

The probability mass function describes a discrete random variable by assigning
probability to each individual value. There is another fundamental way to encode
the same information: instead of asking for the probability of exactly one value,
we ask for the probability that the random variable is at most a given threshold.

\begin{definitionbox}
Let \(X\) be a random variable. Its \textbf{cumulative distribution function}, or
\textbf{CDF}, is the function
\[
F_X(t)=P(X\leq t),
\qquad t\in\mathbb{R}.
\]
\end{definitionbox}

The CDF accumulates probability from the left. For each real number \(t\), it
answers the question:
\[
\text{how much probability has been accumulated up to }t?
\]

\subsection*{The CDF in the discrete case}

If \(X\) is discrete with probability mass function \(p_X\), then
\[
F_X(t)=\sum_{x\leq t}p_X(x),
\]
where the sum is taken over support values \(x\) not exceeding \(t\). Thus the
CDF is obtained by adding the masses at all values to the left of \(t\).

\begin{remarkbox}
The PMF answers ``what is the probability of this exact value?'' The CDF answers
``what is the probability of this value or anything smaller?''
\end{remarkbox}

\subsection*{Example: one fair die}

Let \(X\) be the result of one fair die roll. Then
\[
P(X=k)=\frac16,\qquad k=1,2,3,4,5,6.
\]
The CDF is
\[
F_X(t)=P(X\leq t).
\]
For example,
\[
F_X(0.5)=0,
\]
because no die result is at most \(0.5\). Also,
\[
F_X(1)=\frac16,\qquad F_X(2)=\frac26,\qquad F_X(3.7)=\frac36.
\]
The value \(F_X(3.7)\) equals \(3/6\) because the die results not exceeding
\(3.7\) are \(1,2,3\).

A complete description is
\[
F_X(t)=
\begin{cases}
0, & t<1,\\
\frac16, & 1\leq t<2,\\
\frac26, & 2\leq t<3,\\
\frac36, & 3\leq t<4,\\
\frac46, & 4\leq t<5,\\
\frac56, & 5\leq t<6,\\
1, & t\geq 6.
\end{cases}
\]
This is a step function. It jumps at the possible values of \(X\).

\subsection*{Jumps and probability masses}

For a discrete random variable, the size of the jump of the CDF at a value \(x\)
is exactly the probability mass at that value. If the support values are ordered
as
\[
x_1<x_2<x_3<\cdots,
\]
then
\[
p_X(x_1)=F_X(x_1),
\]
and for \(k\geq 2\),
\[
p_X(x_k)=F_X(x_k)-F_X(x_{k-1})
\]
when there are no support values between \(x_{k-1}\) and \(x_k\).

This means that the PMF and CDF contain the same information in the discrete
case. The PMF stores the masses directly. The CDF stores their cumulative sums.

\subsection*{Example: indicator random variable}

Let \(I_A\) be the indicator of an event \(A\), and write
\[
p=P(A).
\]
Then
\[
P(I_A=1)=p,\qquad P(I_A=0)=1-p.
\]
The CDF is
\[
F_{I_A}(t)=
\begin{cases}
0, & t<0,\\
1-p, & 0\leq t<1,\\
1, & t\geq 1.
\end{cases}
\]
There is a jump of size \(1-p\) at \(0\) and a jump of size \(p\) at \(1\).

\subsection*{Example: sum of two dice}

Let \(S\) be the sum of two fair dice. We know that
\[
\operatorname{supp}(S)=\{2,3,\ldots,12\}.
\]
For example,
\[
F_S(6)=P(S\leq 6).
\]
The sums at most \(6\) are \(2,3,4,5,6\). Their counts among the \(36\) ordered
pairs are
\[
1,2,3,4,5.
\]
Therefore
\[
F_S(6)=\frac{1+2+3+4+5}{36}=\frac{15}{36}=\frac5{12}.
\]
Similarly,
\[
F_S(7)=\frac{1+2+3+4+5+6}{36}=\frac{21}{36}=\frac7{12}.
\]
The difference is
\[
F_S(7)-F_S(6)=\frac6{36}=P(S=7).
\]

\subsection*{Basic properties of a CDF}

Every CDF has three fundamental properties.

First, it is non-decreasing. If \(s<t\), then the event \(\{X\leq s\}\) is
contained in \(\{X\leq t\}\), so
\[
F_X(s)\leq F_X(t).
\]

Second, it approaches \(0\) far to the left and \(1\) far to the right. In finite
discrete examples, this is visible immediately: below the smallest support value
no mass has appeared, and above the largest support value all mass has appeared.

Third, in the discrete case, the CDF is constant between support values and jumps
at support values.

\subsection*{Why use the CDF?}

The CDF is useful because many probability questions are cumulative:
\[
P(X\leq t),\qquad P(X>t),\qquad P(a<X\leq b).
\]
Once the CDF is known, these can be computed systematically. For example,
\[
P(X>t)=1-F_X(t).
\]
If \(a<b\), then
\[
P(a<X\leq b)=F_X(b)-F_X(a).
\]
For discrete variables, endpoint conventions must be read carefully, but the CDF
keeps the logic organized.

\begin{theorembox}
For a discrete random variable, the CDF is the accumulated PMF:
\[
F_X(t)=\sum_{x\leq t}p_X(x).
\]
It describes the same probability law as the PMF, but from the cumulative point
of view.
\end{theorembox}
